\section{Pre-registration and deviations}
\label{app:prereg}

The protocol was written as four frozen phase documents (diagnostic,
characterisation, decomposition, evaluation) plus a research statement, before
any data were collected; a separate progress tracker was the only mutable
record, so no pre-registered bar could be edited after the fact. The
registration is internal rather than platform-hosted, and how far it can be
audited depends on the phase. The four phase documents are plain files in the
released repository, present in its initial commit and thereafter edited
exactly once, to remove a process line (no quantitative content); the
research statement gained one status note and is otherwise unchanged. That
initial commit is dated 22~June 2026, the day Gate~1 was decided, and the
tracker committed with it already records every Phase-1 step as done and
Gate~1 as passed (the Phase-1 run records are dated 20--22~June). The public
git history therefore shows the Phase 2--4 documents --- the stationarity
test and its decision matrix, the Gate-3 criteria, the causal thresholds ---
fixed before the analyses they govern, the first of which was committed later
that day. For Phase~1 (H1, its two nulls, the Tuned-Lens gate and the Gate-1
questions) it shows nothing either way, and that those documents preceded the
data rests on our own record. Commit dates are in any case set by the
committer, not by a third party. The gates:

\begin{itemize}
  \item \textbf{Gate 1} (H1 supported above null?) --- \emph{PASS}, all six
    pre-registered questions answered yes; routed to Phase 2.
  \item \textbf{Gate 2} (H2 stationarity?) --- \emph{FAIL (pervasive drift)};
    routed, per the pre-registered decision matrix, to the state-space branch
    with free per-layer operators.
  \item \textbf{Gate 3} (reconstruction competitive at matched sparsity, and
    seed-stable features?) --- \emph{GO}: reconstruction clause met with
    margin; the seed-stability clause initially failed (0.45 < 0.7) and was
    recorded as a limitation per the plan's ``reconstruction fine / seed
    unstable'' branch, later traced to the source dictionary
    (\cref{sec:evaluation:stability}). The clause was evaluated on refits of
    the SSM over a fixed dictionary (deviation~8 below).
  \item \textbf{Causal existence proof} --- \emph{NULL} at all four
    pre-registered thresholds; reported at the plan's designated honest
    landing place, with no threshold adjusted.
\end{itemize}

\noindent Deviations from the frozen protocol, as completely as we can
enumerate them:

\begin{enumerate}
  \item \emph{Dataset thresholds} (plan-sanctioned tuning direction):
    CounterFact confidence threshold $0.3 \to 0.1$ and per-class floor
    $500 \to 300$, because the primary model knows only ${\sim}2$k facts;
    yield 1{,}555 examples / 3 classes instead of the aspirational
    5--10k / $\ge$4 classes (whose ${\ge}1{,}000$-per-class \emph{target} is
    the figure \cref{sec:limitations} cites; the 500 was the per-class
    inclusion \emph{floor}). Relation-class scope became Limitation~3 by an
    explicit recorded decision (option (b) of the scale-up fork) after the
    cross-model replications made class breadth non-load-bearing.
  \item \emph{Primary checkpoint}: Pythia-410m $\to$ Pythia-410m-\emph{deduped}
    (the public Tuned Lens required for the pre-registered gate baseline
    exists only for deduped).
  \item \emph{Stationarity-test order cap} (in-plan): $k=5$ from the plan's
    early-band feasibility rule; working rank promoted $128 \to 256$ by the
    plan's own $r$-vs-$2r$ spot-check.
  \item \emph{Cross-model targets}: Gemma-2-2B / Llama-3.1-8B do not fit the
    3\,GB study GPU; replaced by the five replication models of
    \cref{tab:crossmodel} that do (with the primary: six models, three
    families; Limitation~2). The plan's Tuned-Lens gate was
    waived-and-reported on replication models without public lenses.
  \item \emph{Statistic repair} (post-hoc, both statistics archived): the
    depth-series gap statistic changed from max-over-$\lambda$ matched-$L_0$
    interpolation to the raw same-$\lambda$ knee gap after per-seed forensics
    showed the former pairs against a saturated baseline tail; the knee
    statistic is the conservative one.
  \item \emph{Post-hoc analyses} beyond the protocol. None of the following
    was specified in the frozen documents, whether or not the text repeats
    the label where the analysis appears: the T1--T5 pressure battery (T4 is
    its untrained-network AR control) and the untrained reconstruction
    analogue; the depth-controlled lag sweep and the common-$k$
    re-derivations; the sublayer split; the object-disjoint split, the
    generic-text replication and the Freedman--Lane calibration;
    parameter-matched baselines; the depth-scaling series; the rank,
    parameter-budget and sample-efficiency axes of the CLT comparison and
    its validation-based reselection; the trained L1 and TopK dictionaries
    and what is derived from them (the refit-stability diagnosis, the gauge
    analysis, the L1 re-run of the causal test); the control ladder of the
    causal section (amplification, co-firing groups, count-, mass- and
    norm-matched random controls, off-target specificity); the operator and
    state-path analysis; the state probe; the source-prefix sweep beyond its
    registered embedding-only point; the paired bootstrap intervals; and the
    dictionary-level seed check of \cref{app:dictseeds}.
  \item \emph{Two self-corrections}, made after the draft was written and
    reported in the open because both had propagated into headline sentences.
    (a) The fitted state operator's spectral radius was described as
    gauge-invariant; in a depth-indexed SSM it is not, and the mechanism claim
    now rests on the invariant state-path decay
    (\cref{sec:evaluation:operator}). (b) The pooled AR($k$) sweep scores a
    depth set that shrinks with $k$, so ``best over $k$'' silently changed the
    evaluation set; the affected contrasts are re-derived at a common $k$ in
    \cref{tab:commondepth}, and the pre-registered headline is reported as the
    four-depth number it is. Neither correction changes a gate outcome.
  \item \emph{Seed-stability check as implemented} (identified in the v1.1
    revision; \cref{app:changes}). The plan's Phase-3 sanity step asked for
    the whole pipeline to be re-run under three seeds ``for the SAE init and
    the optimiser'', scoring features by the maximum cosine of their decoder
    directions across seeds. The sources are frozen public SAEs, so the check
    was implemented on the SSM alone: it is refit from three initialisations
    over one fixed dictionary and features are paired by index
    (\cref{app:stats}). Version 1.0 neither listed this narrowing nor
    described the statistic accurately; within the pre-registered analyses
    no dictionary was retrained under a second seed (the post-hoc check of
    \cref{app:dictseeds} does so), and the Gate-3 stability clause should be
    read accordingly.
  \item \emph{Baselines as run} (the plan's Phase-4 baseline step; this item
    and the next were identified when the paper was re-read against the
    protocol after release). The protocol specified four baselines.
    (A)~A per-layer SAE with its implied update
    $\hat v_{\ell+1} - \hat v_\ell$: run; it is the $-0.223$ of
    \cref{sec:evaluation:baselines}, which \cref{fig:pareto} labels
    ``Baseline~B''. What that figure labels ``Baseline~A'', the per-layer
    dictionary trained directly on $\Dl$, is the comparator we used at
    Gate~3 and a stronger baseline than the plan's. (B)~A crosscoder or
    shared-dictionary
    MLSAE, to be replicated on the same activations where no public weights
    could be used: \emph{not run}. Neither was fit or evaluated (training a
    full-width crosscoder exceeds the study GPU); our MLSAE-derived baseline
    uses that release's independent single-layer SAEs
    (\cref{sec:decomposition:sources}), not a shared dictionary. (C)~The
    Tuned Lens, as its per-layer divergence from the final-layer output
    distribution: we report its implied-update $\Rsq$ instead, the statistic
    the Gate-1 baseline already used. (D)~A CLT in two variants, native
    MLP-write and residual-write, from public weights or a scoped-down
    replication of the published recipe: only the residual-write variant was
    fit, as free decoders over frozen shared sources and not as a natively
    trained CLT (\cref{sec:evaluation:clt}). For the native variant we took
    the plan's fallback and record it as a limitation
    (\cref{sec:limitations}).
  \item \emph{Causal-test candidate selection} (the plan's Phase-4 causal
    step). The plan selected candidates in two stages: features \emph{novel}
    against every non-CLT baseline under two null-calibrated thresholds,
    then filtered for seed and dataset stability. The novelty stage was not
    applied, because the SSM sources from the same atoms as the per-layer
    baselines and no feature can be novel by activation identity
    (Limitation~1, \cref{sec:limitations}). Candidates were chosen by refit
    stability, purity and theme alone.
\end{enumerate}

\section{Operating points and evaluation protocols}
\label{app:operating}

All held-out $\Rsq$ values are against the per-depth-mean baseline on the
same 311-example test split. The apparent spread of ``the SSM's $\Rsq$''
across sections is protocol, not inconsistency:

\begin{table}[t]
  \centering
  \footnotesize
  \caption{Every SSM (and matched-baseline) operating point quoted in this
  paper; ``baseline'' rows are the state-ablated twin
  (\code{use\_state=False}). Gaps are differenced at full precision and then
  rounded, so a row pair is quoted to as many decimals as its gap needs.}
  \label{tab:operating}
  \setlength{\tabcolsep}{4pt}
  \begin{tabular}{llll}
    \toprule
    Value & Protocol & Setting & Used for \\
    \midrule
    0.880 & train fit & $\lambda_{L1}{=}1$, $r{=}d_z{=}32$ & convergence/sanity \\
    0.867 & LASSO held-out & $\lambda_{L1}{=}0.3$ & Pareto top end (\cref{fig:pareto}) \\
    0.857 & LASSO held-out & $\lambda_{L1}{=}1$ (knee) & headline gap $+0.049$ \\
    0.698 & LASSO held-out & $\lambda_{L1}{=}30$ & Pareto bottom end (\cref{fig:pareto}) \\
    0.849 & LASSO held-out & rank-64 baseline, $\lambda_{L1}{=}1$ & param-matched gap $+0.008$ \\
    0.816 & LASSO held-out & baseline, $\lambda_{L1}{=}0.3$ & baseline Pareto top end \\
    0.808 & LASSO held-out & baseline, $\lambda_{L1}{=}1$ (knee) & knee baseline of the $+0.049$ gap \\
    0.653 & LASSO held-out & baseline, $\lambda_{L1}{=}30$ & baseline Pareto bottom end \\
    0.79335 & frozen-source & SSM rank 32, shared grid $\lambda_{\mathrm{reg}}{=}1$ & CLT depth list, $N$ series \\
    0.80383 & frozen-source & CLT rank 32, shared grid & the $-0.0105$ gap \\
    0.7975 / 0.8117 & frozen-source & SSM / CLT rank 32, extended grid & the $-0.0142$ gap (axis 1) \\
    0.7879 / 0.7969 & frozen-source & SSM / CLT rank 32, validation-selected & the $-0.0090$ gap ($N{=}1{,}057$) \\
    0.8155 / 0.843 & frozen-source & SSM rank 64 / CLT rank 64 & budget axis; the $-0.027$ gap \\
    0.7531 / 0.7941 & frozen-source & SSM / CLT rank 16, shared grid & the $-0.041$ rank-16 gap \\
    0.770$^{\ast}$ & frozen-source & rank 32, $\ell_{\max}{=}23$ & prefix sweep (\cref{fig:prefix}) \\
    \bottomrule
  \end{tabular}
\end{table}

The LASSO protocol re-infers test codes against test $\Delta$ with frozen
operators (reconstruction at matched sparsity; identical for both arms of
every gap). The frozen-source protocol fits decoders on fixed cached codes
and isolates decoder parameterisation; its levels sit uniformly lower. Gaps
are only ever computed within a protocol. $^{\ast}$The prefix sweep has its
own fit recipe (one regularisation setting and 250 Adam epochs, against 180
epochs and a $\lambda_{\mathrm{reg}}$ grid for the CLT-family rows), which is
why its full-source SSM ($0.770$) sits below the $0.793$ of the shared-grid
row at the same rank.

\section{Statistical details}
\label{app:stats}

\paragraph{Granger-style nested-F.}
At each (depth $\ell$, lag $j$) cell we compare nested AR models (lags
$1..j{-}1$ vs $1..j$, per-depth intercepts included) on a top-32 PCA
projection of $\Delta$ (full-$D$ Wilks is singular at $D \gg N$), using
Wilks' $\Lambda$ with Rao's F approximation. Bonferroni corrects across the
112 feasible cells ($\alpha = 4.5\times10^{-4}$ per cell); the feasibility
count follows the plan's rule ($\ell \ge j$), i.e.\ $\sum_j (L - j)$ over the
lag grid. With the full grid $\{1,2,3,5,10,15,20\}$ this gives 112 cells at
$L{=}24$ and 196 at $L{=}36$; the 12-layer replications
(\cref{tab:crossmodel}) cap the grid at $\{1,2,3,5,8,10\}$ --- lags 15 and
20 are infeasible at $L{=}12$, with lag 8 substituted to keep six grid
points --- giving 43 cells, with Bonferroni taken
over the per-model feasible count.

\paragraph{Permutation calibration.}
Because the rogue high-variance dimensions make errors heavy-tailed, the F
thresholds were checked with a Freedman--Lane permutation scheme
\citep{freedman1983nonstochastic}: fit the reduced model, permute its
residuals, recompute F (999 permutations), at 12 spot cells spanning mid/late
depth and lags 1--20. All 12 observed F statistics sit at the permutation
floor ($p \le 10^{-3}$), and the permutation null's tail quantiles match the
parametric $F(\mathrm{df}_1,\mathrm{df}_2)$ within 1\% at the 95th/99th
percentiles --- the Bonferroni cut behaves as advertised under
exchangeability. The scrambled-null grids (0/112 significant, both stores)
are the complementary full-grid empirical calibration.

\paragraph{Paired bootstrap CIs.}
Every headline gap is a paired bootstrap over the 311 test examples (10k
resamples) on per-example sums of squares, computed from saved per-example
arrays --- pure post-processing, no refitting. Seed-level uncertainty, where
reported (\cref{fig:depthscaling}), is mean$\pm$sd over three full refits.

\paragraph{Selection, and the AR depth-set confound.}
AR best-($k$, rank) maximises held-out $\Rsq$ over 28 symmetric
real-vs-null evaluations; the real$-$null contrast is the claim. But a pooled
AR($k$) fit scores only depths $\ell \ge k$, so the seven $k$ values are
evaluated on seven different depth sets (23, 22, 21, 19, 14, 9 and 4 target
depths at $L{=}24$). Selecting $k$ by held-out $\Rsq$ therefore partly selects
the easiest depth subset, and the pre-registered headline ($\Rsq = 0.65$ at
$k{=}20$) is an average over $\ell \in [20,23]$, not over the stack. Two
consequences we correct rather than hide:

\begin{itemize}
  \item \emph{Isolating the lag effect.} Re-running the $k$ sweep with the
    depth set held fixed at $\ell \in [20,23]$ separates lag order from depth
    subset, and agrees with the per-depth verdict of \cref{sec:phase2:kstar}:
    order, not reach, is what reconstruction sees.
  \item \emph{Like-for-like contrasts.} The T1 $m$ sweep and the T5 split are
    re-run at fixed $k{=}1$ (all feasible depths) on every store the pressure
    battery uses. \Cref{tab:commondepth} is that table; all trained-vs-untrained
    and factual-vs-generic comparisons in \cref{sec:phase1:pressure} are taken
    from it. The pre-registered headline numbers are reported unchanged
    alongside. One conclusion changed: at a common $k$, GPT-2-small's
    tautology-free AR structure is about half learned, not ${\sim}5\%$, so the
    ``architectural'' reading of T4 is specific to \pythia{}.
\end{itemize}

\begin{table}[t]
  \centering
  \small
  \caption{The T1/T5 pressure battery re-run at a \emph{common} lag order
  ($k{=}1$, every feasible depth), so trained/untrained and factual/generic are
  scored on the same targets. $\Rsq_{m=0}$ is the raw pooled AR(1) level;
  $\Rsq_{m=50}$ strips the top-50 persistent state directions (T1); the last two
  columns are T5's exact orthogonal split of $\Rsq_{m=0}$ at $m{=}50$
  (tautological $+$ genuine $=$ raw). The split is exact at full precision;
  columns are rounded independently, so a displayed pair can miss the
  displayed raw level by one unit in the last digit (first row:
  $0.272 + 0.248$ vs $0.521$). Scrambled nulls are $\le 7\times10^{-3}$
  in every row. Run \code{r13-common-depth-pressure}.}
  \label{tab:commondepth}
  \begin{tabular}{lccccc}
    \toprule
    Store & depths & $\Rsq_{m=0}$ & $\Rsq_{m=50}$ (T1) &
    T5 taut. & T5 genuine \\
    \midrule
    \pythia{} factual, trained    & 23 & 0.521 & \textbf{0.436} & 0.272 & 0.248 \\
    \pythia{} factual, untrained  & 23 & 0.543 & \textbf{0.503} & 0.116 & 0.427 \\
    \pythia{} generic, trained    & 23 & 0.882 & \textbf{0.069} & 0.879 & 0.003 \\
    \pythia{} generic, untrained  & 23 & 0.063 & \textbf{0.061} & 0.007 & 0.056 \\
    GPT-2-small factual, trained   & 11 & 0.626 & \textbf{0.439} & 0.490 & 0.136 \\
    GPT-2-small factual, untrained & 11 & 0.246 & \textbf{0.233} & 0.046 & 0.200 \\
    \bottomrule
  \end{tabular}
\end{table}

\noindent Three readings follow. (i) On \pythia{} the untrained network's
persistence-stripped $\Rsq$ \emph{exceeds} the trained network's
($0.503 > 0.436$), so T4's architectural verdict stands on the primary model.
(ii) On GPT-2-small it does not ($0.439$ trained vs $0.233$ untrained; learned
share $47\%$). (iii) Generic text, compared at the same $k$, retains about a
sixth of the tautology-free structure of factual text ($0.069$ vs
$0.436$) --- the raw generic level of $0.88$ is $99.7\%$ drift.

The CLT
family's $\lambda_{\mathrm{reg}}$ was originally selected on test (house
convention, interior peaks); \cref{sec:evaluation:clt} re-derives the one
consequential CI with validation-based selection, where every arm picks an
interior $\lambda^{*}{=}3$ and the conclusion is unchanged (test selection had
mildly flattered the CLT, not the SSM). The four CLT gap numbers in
\cref{sec:evaluation:clt} map to protocols as follows: $-0.0105$ is the shared
grid ($\lambda_{\mathrm{reg}}{=}1$, the grid every cross-model point uses);
$-0.0142$ is the extended grid, best-of-grid per arm; $-0.0090$ is
validation-selected on a 1{,}057-example train carve-out; $-0.011$ in the
depth list is the shared-grid primary-model point, i.e.\ $-0.0105$ rounded.

\paragraph{Statistics defined once.}
\emph{match@0.7} is a statistic of the SSM, not of the dictionary. The SSM is
refit from different random initialisations over one fixed dictionary (three
refits, so three pairs, averaged). For each of the top-$n$ features by
activation mass in the first fit ($n = 100$ at Gate 3, $n = 200$ in the gauge
analysis), the feature's write direction $D_\ell[:, j]$ at its home layer is
compared between two refits, pairing features \emph{by index}: every refit
shares the same frozen dictionary, so feature $j$ is the same atom in each
and there is no permutation to search over. match@0.7 is the fraction of
those features with cosine $\ge 0.7$. It is silent on whether two
independently trained dictionaries would contain the same atoms;
\cref{app:dictseeds} measures that separately.
\emph{Purity} is the largest relation class's share of a feature's top-100
activating examples. \emph{Retention ratio} is a cross-condition $\Rsq$
divided by the target condition's within-condition $\Rsq$; the tautology,
being common to both, is largely normalised out (it enters additively, so
the cancellation is approximate, not exact). \emph{Knee gap} is the raw same-$\lambda$
SSM$-$baseline gap at $\lambda_{L1}{=}1$.

\section{Cross-model inventory}
\label{app:crossmodel}

\Cref{tab:zoo} lists the models, datasets, and source dictionaries.

\begin{table}[t]
  \centering
  \footnotesize
  \caption{Models, datasets, and source dictionaries. Confidence-filtered
  CounterFact sizes (70m: top-1{,}000 candidates by confidence --- the
  optimiser-debug lane, excluded from the H1/H2 replication claim); SAE
  column gives the public suite used as frozen sources (``---'' = diagnostics
  only). The stability column is match@0.7 of the SSM's per-feature write
  directions across three refits of the SSM over the listed dictionary (the
  dictionary is fixed; \cref{app:stats}); where a
  row names two dictionaries it lists them in the same order, and ``---'' means
  not measured (the SSMs over the GPT-2-small and Qwen2 trained-L1
  dictionaries were fit once each).}
  \label{tab:zoo}
  \begin{tabular}{llllll}
    \toprule
    Model & $L$ & $D$ & $N$ & SAE source (type) & refit stab. \\
    \midrule
    Pythia-70m-deduped   & 6  & 512  & 1{,}000 & EleutherAI 32k (TopK $k{=}16$) & 0.99 \\
    Pythia-160m-deduped  & 12 & 768  & 767     & MLSAE x64-k32 (TopK) & 0.38 \\
    \pythia{}            & 24 & 1024 & 1{,}555 & MLSAE x64-k32 (TopK) / trained L1 & 0.45 / 1.00 \\
    GPT-2-small          & 12 & 768  & 1{,}358 & res-jb (L1) / trained L1 & 0.95 / --- \\
    GPT-2-medium         & 24 & 1024 & 1{,}795 & --- & --- \\
    GPT-2-large          & 36 & 1280 & 1{,}924 & --- & --- \\
    Qwen2-0.5B           & 24 & 896  & 2{,}540 & \citet{chanin2024qwensaes} (L1, 12L prefix) / trained L1 & 0.81 / --- \\
    \bottomrule
  \end{tabular}
\end{table}

Its stability column is the cross-model half of the
\cref{sec:evaluation:stability} finding: over L1 sources the SSM's write
directions are reproduced across refits on two additional families, over
TopK sources they are not on two Pythia models, with dataset size ruled out by
subsampling GPT-2 to the smallest $N$ (match 0.95 unchanged). The debug
lane's EleutherAI TopK suite is nonetheless stable (0.99): the instability
is not universal over TopK sources --- every TopK source tested at $\ge$12
layers was unstable and every L1 source stable --- which is why the
controlled within-model contrast of \cref{sec:evaluation:stability}, not the
zoo, carries the mechanism claim.

\section{Reproducibility}
\label{app:repro}

All experiments ran on a single consumer GPU (NVIDIA GTX 1050, 3\,GB, Pascal)
with a pinned environment (Python 3.12; \code{torch} 2.5.1+cu124;
\code{transformers} $<5$; TransformerLens 2.x); the constraint is documented
because it shaped the design (activation caching, alive-set compression,
fp16 loading for the larger models, and the rank-128 CLT being infeasible).
Every run logs its resolved config, git commit, seed, and library versions
before doing any work; those per-run logs stay on the study machine, and the
released record of a run is its tables and, where one was written, its memo.
Numeric intermediates (CSV/JSON tables) are tracked in
git; heavy artefacts (activation stores, checkpoints) are regenerable from
scripts. A single entry point --- \code{scripts/reproduce.sh} --- checks the
environment, runs the test suite (231 tests), asserts dataset sizes, and
regenerates all 16 paper figures byte-identically from the tracked tables;
\code{--toy-pipeline} additionally re-runs the full Pythia-70m lane end to
end. Code, protocol documents, and the progress tracker are released with
the paper at \url{https://github.com/aljazfrancic/delta-stream}, archived on Zenodo
(\href{https://doi.org/10.5281/zenodo.24036211}{doi:10.5281/zenodo.24036211}); code under the MIT licence, the
paper, protocol, figures, and run records under CC BY 4.0.

Software: PyTorch \citep{paszke2019pytorch}, TransformerLens
\citep{nanda2022transformerlens}, scikit-learn \citep{pedregosa2011sklearn},
statsmodels \citep{seabold2010statsmodels}, and the SAE suites cited in
\cref{sec:decomposition:sources}.

\section{Dictionary-level seed stability (post hoc)}
\label{app:dictseeds}

This check is post hoc, added after the analyses above were complete. Its
design, primary statistic, nulls,
tolerances and verdict criteria were committed, before it was run, to the
run record \nolinkurl{runs/v11-seed-stability/memo.md}. On the primary model we retrained
the L1 and TopK dictionaries of \cref{sec:decomposition:sources} at seeds 0,
1 and 2 with the recorded recipe. Retrained seed~0 reproduces the recorded
dictionaries exactly (alive counts and reconstruction $\Rsq$ on all 23
trained layers), and SSM refits over it reproduce the recorded run tables
byte for byte.

For each pair of seeds and each layer we compare the decoder columns of the
alive latents. \emph{Best-match}: each latent's largest \emph{signed} cosine
against the other seed's alive latents, taken in both directions (codes are
non-negative, so $d$ and $-d$ are different features). The primary
statistic, fixed in advance, is the fraction of the top-100 latents by
activation mass whose best match is $\ge 0.7$, averaged over layers and then
over the three seed pairs; ``stable'' means $\ge 0.7$, the Gate-3 bar.
\emph{PW-MCC} is the mean absolute cosine under the optimal one-to-one
assignment \citep{song2026consistency}, restricted to alive latents. The two
nulls are a random isotropic dictionary of matched size and the other seed's
dictionary at the layer half the stack away. \emph{Init retention} is an
alive latent's cosine with its own initial decoder column.

\begin{table}[t]
  \centering
  \small
  \caption{Dictionary-level seed stability on \pythia{} (three seeds per
  mechanism, 23 layers; post hoc). The last row is the SSM-level statistic
  of \cref{sec:evaluation:stability}, refit over each dictionary seed.}
  \label{tab:dictseeds}
  \begin{tabular}{lll}
    \toprule
    Statistic & L1 & TopK \\
    \midrule
    match@0.7, top-100 latents (primary) & 0.254 & 0.035 \\
    \quad per seed pair & 0.251 / 0.255 / 0.255 & 0.035 / 0.040 / 0.030 \\
    \quad range over layers & 0.08--0.42 & 0.00--0.15 \\
    \quad random / cross-layer null & 0.000 / 0.000 & 0.000 / 0.000 \\
    match@0.7, all alive latents & 0.033 & 0.037 \\
    mean best-match cosine, top-100 (random null) & 0.52 (0.10) & 0.31 (0.09) \\
    PW-MCC, alive latents (random null) & 0.17 (0.10) & 0.34 (0.09) \\
    alive latents with init cosine $\ge 0.7$ & 88\% & 32\% \\
    \quad among the top-100 & 31\% & 56\% \\
    \midrule
    SSM refit match@0.7, dictionary seeds 0 / 1 / 2 & 1.00 / 1.00 / 1.00 & 0.08 / 0.06 / 0.08 \\
    \bottomrule
  \end{tabular}
\end{table}

Neither mechanism reaches the bar and both nulls are at zero, so the
pre-committed verdict is that neither gives reproducible dictionaries in this
regime. The real cosines sit well above their nulls, so the seeds do share
some structure. The two mechanisms are ranked differently by different
statistics, and the L1 dictionaries are largely initialisation-dominated:
88\% of alive L1 latents keep cosine $\ge 0.7$ with their own seed-specific
starting column, which no other seed shares. Both observations describe this
regime --- more latents than training vectors, full-batch Adam for 1{,}000
steps, no auxiliary loss --- not SAEs trained at scale, where
\citet{paulo2026seeds} find L1 SAEs more seed-stable than TopK ones and
\citet{song2026consistency} report PW-MCC around 0.80 for TopK. The SSM-level
refit contrast holds on every draw.

\section{Changes since version 1.0}
\label{app:changes}

Version 1.1 changes no pre-registered analysis and no recorded result. It
adds one post-hoc experiment, which retrains the primary model's trained
dictionaries at further seeds, and it corrects statements that version~1.0
worded more strongly than its own run tables support. Everything below is
post hoc. The first seven items concern how the stability finding is worded,
what the paper is compared against, and the new experiment; the last five
are corrections found by re-reading the paper against its tables, code,
protocol and repository history, the final one from a second such pass.

\begin{itemize}
  \item \emph{What ``seed'' means in the stability finding.} Version 1.0
    said that TopK dictionaries make single-feature directions
    ``seed-unidentifiable'', that matched L1 dictionaries are ``perfectly
    stable'', and that TopK features should be checked to ``survive a
    re-seed''. The statistic behind those sentences re-initialises the SSM
    over one fixed dictionary and pairs features by index; no dictionary was
    retrained. Every statement of the finding --- abstract, contributions,
    \cref{sec:related,sec:decomposition:sources,sec:evaluation:stability,sec:discussion},
    the captions of \cref{fig:depthscaling,fig:stability}, and
    \cref{tab:zoo} --- now says what is refit, and the practitioner caution
    is restricted to models fit downstream of a dictionary. The numbers are
    unchanged.
  \item \emph{Definition corrected.} \Cref{app:stats} defined match@0.7 by a
    feature's ``best match in a re-seeded fit''. The code pairs features by
    index; the text was wrong and is fixed.
  \item \emph{Deviation added.} The SSM-only implementation of the plan's
    seed-stability step is now listed in \cref{app:prereg} (deviation~8).
  \item \emph{Regime stated.} The trained dictionaries have more latents
    than training vectors per layer, are fit full-batch without an auxiliary
    loss, and were one draw per mechanism in version~1.0
    (\cref{sec:decomposition:sources,sec:limitations}).
  \item \emph{New post-hoc experiment.} Two further dictionary seeds per
    mechanism on the primary model (\cref{app:dictseeds}). The SSM-level
    refit contrast holds on each draw (L1 1.00, TopK 0.06--0.08); the
    dictionaries themselves do not reproduce across seeds under either
    mechanism in this regime. The abstract, contributions,
    \cref{sec:related,sec:decomposition:sources,sec:evaluation:stability,sec:limitations,sec:discussion}
    and the captions of \cref{fig:depthscaling,fig:stability} say so.
  \item \emph{Figure labels.} \Cref{fig:depthscaling,fig:stability} label
    the statistic ``refit'' stability where version~1.0 said ``seed''; the
    in-panel note of \cref{fig:stability} now records that the L1 result holds
    over three dictionary seeds. No plotted value changed.
  \item \emph{Related work added.} SAE seed consistency
    \citep{paulo2026seeds,song2026consistency}, depth as a dynamical system
    \citep{fernando2025transformer,fernando2026dynamics}, and random-baseline
    sanity checks for SAEs \citep{korznikov2026sanity}.
  \item \emph{Statements corrected against the run tables.}
    (i)~\Cref{sec:evaluation:stability} said the three dictionaries gave
    ``the same reconstruction win ($+0.041$--$+0.047$)''. That range was the
    L1 dictionary's alone, under the interpolated statistic that deviation~5
    retired. The knee gaps are $+0.048$ (L1), $+0.038$ (trained TopK) and
    $+0.049$ (MLSAE), and the trained-TopK fit falls below its baseline at
    the two sparsest points; the section heading and the first finding's
    title now claim the contrast, not a general mechanism.
    (ii)~\Cref{sec:evaluation:prefix} presented $-0.44$ as the
    embedding-only floor. It is an under-fit bias at the final layer, which
    also depresses the three cutoffs below $\ell_{\max}{=}15$; the text and
    the caption of \cref{fig:prefix} say so, and the plotted values are
    unchanged.
    (iii)~\Cref{sec:evaluation:clt,sec:discussion} and the contributions
    stated ``most of the cross-layer value at ${\sim}16\%$ of the
    parameters'' and the sub-$10^3$-example crossover without restricting
    them to the primary model. Both are now scoped, the per-model shares
    ($30$--$85\%$) are given, and the matched-rank gaps at ranks 16 and 64
    are stated.
    (iv)~\Cref{sec:evaluation:probe} read the state's lead over raw-residual
    PCA at $\ell{=}17$ as the state running a layer ahead. The probe's codes
    are fit to each example's whole trajectory, so that comparison is
    confounded; the text, the caption of \cref{fig:probe} and
    \cref{sec:discussion} say so. The contrasts among code-derived arms are
    unaffected.
    (v)~The trained L1 dictionary is no longer called ``clean'' or
    ``cleaner'', and the causal null is described as holding under both
    dictionaries tested, not as dictionary-independent: the new experiment
    shows that dictionary does not reproduce across seeds.
    (vi)~The object-disjoint AR level in \cref{sec:phase1:verdict} now
    carries its depth set and the whole-stack value ($0.411$ against
    $0.521$).
  \item \emph{Deviation list completed.} \Cref{app:prereg} gains
    deviations~9 (baselines as run) and~10 (the causal test's novelty
    filter), deviation~6 lists the post-hoc analyses in full as far as we
    can identify them, and the list no longer calls itself exhaustive. The
    abstract, the introduction and
    \cref{sec:evaluation} mark the causal control ladder, the operator
    analysis, the depth series and the three-axis CLT reading as post hoc.
  \item \emph{Confounds named, checks not run.} Two checks that fit the
    study GPU follow from (ii) and (iv) and are listed as future work in
    \cref{sec:discussion}.
  \item \emph{Small corrections.} A backward reference in
    \cref{sec:phase1:sublayer} said ``below''; ``no control'' in
    \cref{sec:evaluation:causal} is now ``no magnitude-matched control'';
    the alive-latent counts in \cref{sec:decomposition:sources} are given
    for all three seeds; the write-direction subspace overlap
    (${\sim}0.6$) is no longer listed among the stable statistics; the
    summary of \citet{fernando2026dynamics} follows their body text; and
    the test count in \cref{app:repro} is 231, up from 208, for the new
    tooling.
  \item \emph{Second re-read.} Three statements were wrong and six claimed
    more than their evidence; all nine were in version~1.0.
    (i)~\Cref{app:prereg} said the state of the phase documents at every
    analysis step was verifiable from the public git history. That history
    begins after Gate~1, so it confirms the freeze for Phases 2--4 only;
    the appendix, \cref{sec:phase1:setup} and \cref{sec:limitations} now say
    so.
    (ii)~\Cref{sec:phase2:stationarity,sec:phase2:scope} read principal
    angles of 57--62$^\circ$ against a 90$^\circ$ reference and concluded
    that the kernels share structure. For rank-256 subspaces of
    $\mathbb{R}^{1024}$ chance is ${\approx}63^\circ$; the angles are at
    chance and the conclusion is withdrawn.
    (iii)~\Cref{sec:phase2:stationarity} said long-horizon stationarity was
    rechecked on the fitted model and failed. The recheck is defined for the
    branch the study did not take and never ran; stationarity beyond lag 5
    is untested.
    (iv)~The growth of the reconstruction gain with depth holds for the knee
    gap under the LASSO protocol and not for the same contrast under the
    frozen-source protocol, where the six-layer model gains as much as the
    24-layer one. \Cref{sec:evaluation:depth} carries this as a sixth
    caveat; the abstract, the introduction and \cref{sec:evaluation:clt}
    name the protocol.
    (v)~\Cref{sec:evaluation:gate3} called the headline gap
    dictionary-independent. It is the same over two dictionaries, $+0.038$
    over a third, and on GPT-2-small it is $+0.0175$ or $+0.143$ depending
    on the dictionary; the second number was not reported before.
    (vi)~\Cref{sec:evaluation:clt,sec:discussion} and the contributions
    credited the SSM's parameter efficiency against the CLT to the state. A
    per-layer decoder of the same size does as well under that protocol,
    and the text now says so.
    (vii)~The untrained-network shares of the gain ($53\%$, $44\%$) compare
    gaps at equal penalty; the untrained SSM runs 1.3--1.8$\times$ denser
    than its baseline, so they are upper estimates of the architectural
    share (\cref{sec:evaluation:depth}).
    (viii)~T2 is zero by construction, and ``tautology'' is an attribution
    by subspace, not a mechanical identity
    (\cref{sec:phase1:setup,sec:phase1:pressure}); \cref{sec:discussion} no
    longer calls that share definitional.
    (ix)~The state path's short memory is conditional on sources at every
    layer: the prefix sweep's own rows show a refit state carrying signal
    six layers. \Cref{sec:evaluation:operator,sec:evaluation:prefix} no
    longer say that the increments forget or that propagation cannot
    substitute.
    Smaller: three ``order of magnitude'' statements are now ``about six
    times'' ($0.069$ against $0.436$); the probe's scope sentence gives the
    early and mid-depth comparison with the raw stream and the state's
    relation-class accuracy; \cref{tab:operating} explains its prefix-sweep
    row; \cref{tab:zoo} lists the GPT-2-small trained-L1 dictionary;
    \cref{app:repro} says the per-run logs are not part of the release; and
    the title of the middle panel of \cref{fig:clt} no longer calls the CLT
    edge small on every model (no plotted value changed).
\end{itemize}

